\chapter{Solving Linear Systems}
\label{ch:solving-linear-systems}

\section*{Chapter Preview}
\addcontentsline{toc}{section}{Chapter Preview}

The equation $Ax=b$ asks whether $b$ can be built as a linear combination of
the columns of $A$. This chapter classifies what can happen and describes the
shape of the solution set.

\section{Existence}

Let $A\in\R^{m\times n}$ with columns $a_1,\ldots,a_n$. The equation
\[
  Ax=b
\]
means
\[
  x_1a_1+\cdots+x_na_n=b.
\]
Therefore:

\begin{proposition}
$Ax=b$ has a solution if and only if $b\in\Img A$.
\end{proposition}

This is the most important first test. If $b$ is outside the image, the exact
system has no solution. Later, least squares will replace $b$ by its closest
point in the image.

\section{Uniqueness}

Suppose $x_0$ is one solution to $Ax=b$. Then another vector $x$ is also a
solution exactly when
\[
  A(x-x_0)=0.
\]
Thus $x-x_0\in\Ker A$.

\begin{theorem}[Solution set]
If $Ax=b$ has one solution $x_0$, then the full solution set is
\[
  x_0+\Ker A=\{x_0+k:k\in\Ker A\}.
\]
\end{theorem}

So uniqueness is controlled by the kernel:
\[
  \Ker A=\{0\}
  \quad\Longleftrightarrow\quad
  \text{every consistent system }Ax=b\text{ has at most one solution.}
\]

\section{The Four Cases}

The existence question and the uniqueness question are independent.

\[
\begin{array}{c|c|c}
 & b\in\Img A & b\notin\Img A\\
\hline
\Ker A=\{0\}
& \text{one solution}
& \text{no solution}\\
\Ker A\ne\{0\}
& \text{infinitely many solutions }x_0+\Ker A
& \text{no solution}
\end{array}
\]

\begin{warningbox}
The phrase ``the solution'' is only justified after checking uniqueness. In an
underdetermined consistent system, a numerical routine may return one
solution, often the least-norm solution, but the mathematical solution set is
an affine set $x_0+\Ker A$.
\end{warningbox}

\section{Inverses}

For a square matrix $A\in\R^{n\times n}$, the following conditions are
equivalent:
\[
  A^{-1}\text{ exists}
  \quad\Longleftrightarrow\quad
  \Ker A=\{0\}
  \quad\Longleftrightarrow\quad
  \Img A=\R^n
  \quad\Longleftrightarrow\quad
  \Rank A=n.
\]
When these hold, every $b\in\R^n$ has the unique solution $x=A^{-1}b$.

\section{Computation}

In practice, one usually solves linear systems by elimination or by numerical
factorizations, not by explicitly computing $A^{-1}$. For a square full-rank
matrix, \texttt{np.linalg.solve(A, b)} is the standard call. For rectangular
or inconsistent systems, \texttt{np.linalg.lstsq(A, b, rcond=None)} returns a
least-squares solution.

\begin{lstlisting}
import numpy as np

A = np.array([[1., 2.],
              [3., 4.]])
b = np.array([1., 0.])

x = np.linalg.solve(A, b)
print(A @ x)

B = np.array([[1., 1.],
              [1., 2.],
              [1., 3.]])
y = np.array([1., 2., 2.])

x_hat, residuals, rank, s = np.linalg.lstsq(B, y, rcond=None)
print(x_hat)
print(np.linalg.norm(B @ x_hat - y))
\end{lstlisting}

\section*{Chapter Summary}
\addcontentsline{toc}{section}{Chapter Summary}

\begin{itemize}
  \item $Ax=b$ is solvable exactly when $b$ lies in the image of $A$.
  \item If one solution exists, all solutions are $x_0+\Ker A$.
  \item The kernel controls uniqueness.
  \item A square matrix is invertible exactly when its kernel is trivial.
  \item Inconsistent systems lead naturally to least squares.
\end{itemize}

\section*{Where This Appears Later}
\addcontentsline{toc}{section}{Where This Appears Later}

Least squares begins when $b\notin\Img A$. Orthogonal projection explains the
closest solvable right-hand side. The pseudoinverse and SVD later organize all
four cases in one formula.

\begin{labbox}
Lab 4 builds examples for all four cases, parameterizes $x_0+\Ker A$, and
visualizes least-norm solutions.
\end{labbox}

\section*{Exercises}
\addcontentsline{toc}{section}{Exercises}
\subsection*{A. Theory}
\begin{exercise}\exerciseid{6.A1} Explain why coordinates depend on a basis, but the vector being described does not. \end{exercise}
\subsection*{B. By Hand}
\begin{exercise}\exerciseid{6.B1} For $b_1=(1,1)$ and $b_2=(1,-1)$, convert $v=(4,2)$ to $B$-coordinates and back. \end{exercise}
\subsection*{C. Python / NumPy}
\begin{exercise}\exerciseid{6.C1} Write \texttt{change\_coords(B, v)} using \texttt{np.linalg.solve}. Test it on two bases of $\R^2$. \end{exercise}
\subsection*{D. Challenge}
\begin{exercise}\exerciseid{6.D1} Find a basis in which the map $T(x,y)=(x,x+y)$ has a simpler interpretation. \end{exercise}

